\documentclass[../../main]{subfiles}

\ifSubfilesClassLoaded{%
    \zexternaldocument*{../1.Category-Theory/exported/aux/main}%
    \zexternaldocument*{../2.Homological-Algebra/exported/aux/main}%
    \zexternaldocument*{../3.Sheaves/exported/aux/main}%
    \zexternaldocument*{1.Cohomology-of-sheaves}
    \zexternaldocument*{2.de-Rham-cohomology}
    \zexternaldocument*{3.Cech-cohomology}
}{}

\begin{document}
\chapter{Singular cohomology}

\begin{definition}[Standard simplex]\label{def:standard-simplex}
  The \textbf{standard \(n\)-simplex} \(\Delta^{n}\) is the subset of \(\mathbb{R}^{n+1}\) defined by
  \begin{equation*}
    \Delta^{n} \coloneq \left\{ (t_{0}, \ldots, t_{n+1}) \in [0, 1]^{n+1} \mid \sum_{i=0}^{n+1} t_{i} = 1 \right\}
  \end{equation*}
\end{definition}

\begin{definition}[Simplex]\label{def:simplex}
  Let \(X\) be a topological space.
  A \textbf{singular \(n\)-simplex} in \(X\) is a continuous map
  \begin{equation*}
    \sigma: \Delta^{n} \to X
  \end{equation*}
  where \(\Delta^{n}\) is the \hyperref[def:standard-simplex]{standard \(n\)-simplex}.
\end{definition}


\begin{definition}[Singular chains]\label{def:singular-chain}
  The group of \textbf{singular \(n\)-chains} with coefficients in an abelian group \(G\), denoted \(C_{n}(X, G)\), is the free abelian group generated by the set of all \hyperref[def:simplex]{singular \(n\)-simplices} in \(X\).

  We have \textit{boundary} maps \(\delta_{n}: C_{n}^{\operatorname{sing}}(X, G) \to C_{n-1}^{\operatorname{sing}}(X, G)\) defined as
  \begin{equation*}
    \forall \sigma \in C_{n}^{\operatorname{sing}}(X, G) \quad \delta_{n}(\sigma) \coloneq \sum_{i=0}^{n} (-1)^{i} \sigma|_{[v_{0}, \ldots, \hat{v_{i}}, \ldots, v_{n}]}
  \end{equation*}
  A straightforwards calculations shows that \(\delta_{n-1} \circ \delta_{n} = 0\).
  Therefore we have a \textbf{singular chain} \hyperref[def:complex]{complex}
  \begin{center}
    \begin{tikzcd}
      C_{\bullet}^{\operatorname{sing}}(X, G) = \cdots \arrow[r] & C_{n+1}^{\operatorname{sing}} \arrow[r, "\delta_{n+1}"] & C_{n}^{\operatorname{sing}} \arrow[r, "\delta_{n}"] & C_{n-1} \arrow[r, "\delta_{n-1}"] & \cdots \arrow[r, "\delta_{1}"] & C_{0}^{\operatorname{sing}}
    \end{tikzcd}
  \end{center}
\end{definition}

\begin{definition}[Singular cochains]\label{def:singular-cochain}
  The group of \textbf{singular \(n\)-cochains} with coefficients in an abelian group \(G\), denoted \(C_{\operatorname{sing}}^{n}(X, G)\), is the dual of the group of \hyperref[def:singular-chain]{singular \(n\)-chains}:
  \begin{equation*}
    C_{\operatorname{sing}}^{n}(X, G) \coloneq \Hom_{\Ab}(C_{n}^{\operatorname{sing}}(X, G), G)
  \end{equation*}

  We have \textit{coboundary} maps \(d^{n}: C_{\operatorname{sing}}^{n}(X, G) \to C_{\operatorname{sing}}^{n+1}(X, G)\) defined as
  \begin{equation*}
    \forall \phi \in C_{\operatorname{sing}}^{n}(X, G), \sigma \in C_{n+1}^{\operatorname{sing}}(X, G) \quad (d^{n}\phi)(\sigma) \coloneq \phi(\delta^{n}(\sigma)) = \sum_{i=0}^{n+1} (-1)^{i}\phi(\sigma|_{[v_{0}, \ldots, \hat{v_{i}}, \ldots, v_{n+1}]})
  \end{equation*}
  As before, it is easy to check that \(d^{n+1} \circ d^{n} = 0\), so we have a \textbf{singular cochain} \hyperref[def:complex]{complex}
  \begin{center}
    \begin{tikzcd}
      C_{\operatorname{sing}}^{\bullet}(X, G) = C_{\operatorname{sing}}^{0}(X, G) \arrow[r, "d^{0}"] & C_{\operatorname{sing}}^{1}(X, G) \arrow[r, "d^{1}"] & \cdots
    \end{tikzcd}
  \end{center}
\end{definition}

\begin{definition}[Singular cohomology]\label{def:singular-cohomology}
  Let \(X\) be a topological space and \(G\) an abelian group.
  The \textbf{singular \hyperref[def:cohomology]{cohomology}} of \(X\) with coefficients in \(G\) is the \hyperref[def:cohomology]{cohomology} of the \hyperref[def:singular-cochain]{singular cochain complex}:
  \begin{equation*}
    \operatorname{H}_{\operatorname{sing}}^{\bullet}(X, G) \coloneq \cohomology{\bullet}(C_{\operatorname{sing}}^{\bullet}(X, G))
  \end{equation*}
\end{definition}

\begin{lemma}\label{lemma:singular-cohomology-contractible-space}
  Let \(X\) be a contractible space.
  Then for all abelian group \(G\) and \(n > 0\)
  \begin{equation*}
    \operatorname{H}_{\operatorname{sing}}^{n}(X, G) = 0
  \end{equation*}
\end{lemma}

\begin{theorem}\label{thm:small-chains}
  Theorem of small chains
\end{theorem}


\begin{theorem}\label{thm:singular-sheaf-cohomology}
  Let \((X, \sheaf{O}_{X})\) be a \hyperref[def:ringed-space]{ringed space} where \(X\) is locally contractible\footnote{A topological space \(X\) is locally contractible if any open subset \(U \subseteq X\) can be covered by contractible open sets.}.
  Let \(G\) be a commutative ring.

  Then
  \begin{equation*}
    \cohomology{\bullet}_{\operatorname{sing}}(X, G) \cong \cohomology{\bullet}(X, \underline{G})
  \end{equation*}
\end{theorem}
\begin{proof}
  Consider the \hyperref[def:sheaf]{sheaf} \(\sheaf{C}_{\operatorname{sing}}^{q}(-, G)\) associated to the \hyperref[def:presheaf]{presheaf}
  \begin{equation*}
    U \mapsto C_{\operatorname{sing}}^{q}(U, G)
  \end{equation*}
  The differential \(d: C_{\operatorname{sing}}^{q}(U, G) \to C_{\operatorname{sing}}^{q+1}(U, G)\) gives rise to a differential
  \begin{equation*}
    d: \sheaf{C}_{\operatorname{sing}}^{q} \to \sheaf{C}_{\operatorname{sing}}^{q+1}
  \end{equation*}

  This \hyperref[def:complex]{complex} is a \hyperref[def:resolution]{resolution} of \(\underline{G}\).

  Indeed, the complex \(C_{\operatorname{sing}}^{> 0}(U, G)\) is \hyperref[def:exact-complex]{exact} on each contractible open set \(U\) using that the \hyperref[def:cohomology]{cohomology} of the \hyperref[def:complex]{complex} is by definition \hyperref[def:singular-cohomology]{singular cohomology} and \autoref{lemma:singular-cohomology-contractible-space}.

  Since \(X\) is locally contractible every point \(p \in X\) has a neighbourhood basis consisting of contractible open sets.
  The \hyperref[def:stalk-sheaf]{stalk} of the \hyperref[def:complex]{complex} is the \hyperref[def:direct-limit]{direct limit} over this neighbourhoods, so the \hyperref[def:complex]{complex} of \hyperref[def:presheaf]{presheaves} \(C_{\operatorname{sing}}^{\bullet}(-, G)\) is exact at the level of stalks, making the \hyperref[def:complex]{complex} of \hyperref[def:sheaf]{sheaves} \(\sheaf{C}_{\operatorname{sing}}^{\bullet}(-, G)\) is \hyperref[def:exact-complex]{exact}.

  Furthermore since \(X\) is locally pathwise connected we have
  \begin{equation*}
    \ker \left( d^{0}: \sheaf{C}_{\operatorname{sing}}^{0}(-, G) \to \sheaf{C}_{\operatorname{sing}}^{1}(-, G) \right) = \underline{G}
  \end{equation*}

  Note that the \hyperref[def:presheaf]{presheaf} \(C_{\operatorname{sing}}(-, G)\) is \hyperref[def:flasque-sheaf]{flasque}.
  Indeed for any restriction of open sets \(V \subseteq U\) the restricion map is surjective because any \hyperref[def:singular-cochain]{singular cochain} on \(V\) can be extended to one in \(U\) by assigning it to \(0\) to any simplex in \(U\) not contained in \(V\).
  Furthermore we have
  \begin{equation*}
    \sheaf{C}_{\operatorname{sing}}^{q}(U, G) = \quot{C_{\operatorname{sing}}^{q}(U, G)}{C_{\operatorname{sing}}^{q}(U, G)_{0}}
  \end{equation*}
  were \(C_{\operatorname{sing}}^{q}(U, G)_{0}\) is the set of \hyperref[def:singular-cochain]{singular \(q\)-cochains} \(\alpha \in C_{\operatorname{sing}}^{q}(U, G)\) such that there exists a covering \(\mathcal{V}\) of \(U\) with \(\alpha|_{C_{q}(V, G)} = 0\) for all \(V \in \mathcal{V}\).

  Therefore by \autoref{thm:dimension-shifting} we have
  \begin{equation*}
    \cohomology{q}(X, G) = \cohomology{q}(\Gamma(\sheaf{C}_{\operatorname{sing}}^{\bullet}(X, G)))
  \end{equation*}

  It remains to show that the \hyperref[def:complex]{complex}
  \begin{equation*}
    \Gamma(\sheaf{C}_{\operatorname{sing}}^{\bullet}(X, G)) = \quot{C_{\operatorname{sing}}^{\bullet}(X, G)}{C_{\operatorname{sing}}^{\bullet}(X, G)_{0}}
  \end{equation*}
  is \hyperref[def:quasi-isomorphism]{quasi-isomorphic} to the \hyperref[def:complex]{complex} \(C_{\operatorname{sing}}^{\bullet}(X, G)\), which is \autoref{thm:small-chains}.
\end{proof}


\end{document}
